% 我本地环境是 Windows + MiKTeX，使用 VS Code 编写并通过 LaTeX Workshop 编译
% 注意确保 XeLaTeX 完整编译至少两次  目录内容需要第一遍生成 .toc 文件  页码同样依赖第二遍读取lastpage标签
% 本文按 Windows + MiKTeX/TeX Live 配置中文字体；不要使用 pdfLaTeX 编译。
% !LW recipe=latexmk (xelatex)
\documentclass[UTF8,zihao=-4,a4paper,fontset=windows]{ctexart}

%==================== 我的基本信息====================
\newcommand{\StudentName}{Rpc}
\newcommand{\StudentID}{This is my student ID}
\newcommand{\Major}{计算机科学与技术}
\newcommand{\CourseName}{学业基本情况介绍与计算机基本使用}
\newcommand{\ReportDate}{2026年7月18日}
%======================================================


%============用到的一些包================
\usepackage{float}
\usepackage{placeins}
\usepackage{needspace}
\usepackage[a4paper,top=2.5cm,bottom=2.4cm,left=2.6cm,right=2.6cm,headheight=15pt]{geometry}
\usepackage{xcolor}
\usepackage{graphicx}
\usepackage{booktabs}
\usepackage{tabularx}
\usepackage{array}
\usepackage{enumitem}
\usepackage{amsmath,amssymb}
\usepackage{listings}
%===============这里用什么都感觉怪怪的 是用来做示例的  就用demo吧===========
\renewcommand{\lstlistingname}{demo}
\usepackage[most]{tcolorbox}
\usepackage{tikz}
\usetikzlibrary{arrows.meta,positioning,shapes.geometric}
\usepackage{fancyhdr}
\usepackage{lastpage}
\usepackage{hyperref}
\usepackage{microtype}
\usepackage{multicol}

%==========稍稍的美化一下（ps：我只有工业风审美）=====
\definecolor{MainBlue}{HTML}{176B87}
\definecolor{DeepBlue}{HTML}{12344D}
\definecolor{LightBlue}{HTML}{EAF5F8}
\definecolor{Accent}{HTML}{E48B29}
\definecolor{CodeBG}{HTML}{F5F7F9}
\definecolor{Green}{HTML}{39755A}
\definecolor{TextGray}{HTML}{4A5560}

\hypersetup{
  colorlinks=true,
  linkcolor=MainBlue,
  urlcolor=MainBlue,
  citecolor=Green,
  pdftitle={第一次课程学习总结与计算机基础快速入门},
  pdfauthor={\StudentName}
}
\renewcommand{\contentsname}{目录}
\Urlmuskip=0mu plus 1mu

\setlength{\parindent}{2em}
\setlength{\parskip}{0.25em}
\renewcommand{\baselinestretch}{1.35}
\setlist{itemsep=0.25em,topsep=0.35em,leftmargin=2.2em}

\ctexset{
  section={format=\Large\bfseries\color{DeepBlue},number=\chinese{section}、,beforeskip=1.2em,afterskip=0.7em},
  subsection={format=\large\bfseries\color{MainBlue},number=\arabic{section}.\arabic{subsection},beforeskip=1em,afterskip=0.5em},
  subsubsection={format=\normalsize\bfseries\color{TextGray},number=\arabic{section}.\arabic{subsection}.\arabic{subsubsection}}
}

\pagestyle{fancy}
\fancyhf{}
\fancyhead[L]{\small\color{TextGray}\CourseName}
\fancyhead[R]{\small\color{TextGray}第一次课程学习报告}
\fancyfoot[C]{\small 第\thepage{}页，共\pageref{LastPage}页}
\renewcommand{\headrulewidth}{0.4pt}
\renewcommand{\footrulewidth}{0pt}

\newcolumntype{Y}{>{\raggedright\arraybackslash}X}
\newcolumntype{C}[1]{>{\centering\arraybackslash}p{#1}}

\newtcolorbox{summarybox}[1]{
  enhanced,breakable,colback=LightBlue,colframe=MainBlue,
  boxrule=0.7pt,arc=2mm,left=2mm,right=2mm,top=1.5mm,bottom=1.5mm,
  title=#1,fonttitle=\bfseries,coltitle=DeepBlue
}
\newtcolorbox{tipbox}[1]{
  enhanced,breakable,colback=white,colframe=Accent,
  borderline west={2.5pt}{0pt}{Accent},boxrule=0.4pt,arc=1mm,
  title=#1,fonttitle=\bfseries,coltitle=Accent
}

\lstdefinestyle{codestyle}{
  backgroundcolor=\color{CodeBG},basicstyle=\ttfamily\small,
  keywordstyle=\color{MainBlue}\bfseries,stringstyle=\color{Green},
  commentstyle=\color{TextGray}\itshape,numbers=left,numberstyle=\tiny\color{TextGray},
  numbersep=8pt,frame=single,rulecolor=\color{LightBlue},
  breaklines=true,showstringspaces=false,tabsize=4,columns=fullflexible,
  xleftmargin=1em,framexleftmargin=0.7em
}
\lstset{style=codestyle}
\lstdefinelanguage{JavaScript}{
  keywords={async,await,break,case,catch,const,continue,default,delete,do,else,
    export,extends,false,finally,for,function,if,import,in,instanceof,let,new,
    null,return,static,switch,this,throw,true,try,typeof,undefined,var,while},
  sensitive=true,
  comment=[l]{//},morecomment=[s]{/*}{*/},
  morestring=[b]',morestring=[b]"
}

\begin{document}

%-------------------- 封面 --------------------
\begin{titlepage}
  \thispagestyle{empty}
  \begin{tikzpicture}[remember picture,overlay]
    \fill[DeepBlue] (current page.north west) rectangle ([yshift=-3.2cm]current page.north east);
    \fill[MainBlue!18] ([yshift=2.0cm]current page.south west) rectangle (current page.south east);
    \draw[MainBlue,line width=1.2pt] ([xshift=2.5cm,yshift=-3.8cm]current page.north west)--([xshift=-2.5cm,yshift=-3.8cm]current page.north east);
  \end{tikzpicture}
  \vspace*{1.0cm}
  {\color{white}\Large\bfseries 课程学习报告}\par
  \vspace{3.0cm}
  \begin{center}
    {\Huge\bfseries\color{DeepBlue}第一次课程学习总结}\par
    \vspace{0.5cm}
    {\LARGE\color{MainBlue}——计算机基础快速入门与工具实践}\par
    \vspace{1.6cm}
    \begin{tcolorbox}[width=0.82\textwidth,colback=LightBlue,colframe=MainBlue,arc=2mm,boxrule=0.8pt]
      \centering\large
      本报告围绕研究生阶段学习规划、Linux、Markdown、\LaTeX{}以及 Python/C++ 展开，\par
      既总结了课堂内容，也可作为直接使用的个人速查手册。
    \end{tcolorbox}
    \vfill
    \renewcommand{\arraystretch}{1.6}
    \begin{tabular}{>{\bfseries\color{DeepBlue}}r l}
      课程名称：& \CourseName \\
      姓名：& \StudentName \\
      学号：& \StudentID \\
      专业：& \Major \\
      提交日期：& \ReportDate
    \end{tabular}
    \vfill
  \end{center}
\end{titlepage}

\pagenumbering{Roman}
\begin{summarybox}{摘要}
第一次课程从研究生阶段的培养目标出发，介绍了课程学习、科研评价、课题组协作和计算机工具使用等内容。
本文首先梳理课堂要点并形成个人行动计划；随后围绕 Linux、Markdown、\LaTeX{}及编程语言进行系统归纳。
在语言部分，除复习 Python 与 C++ 外，还总结了我较为熟练的 Java、JavaScript、SQL 与 Web 开发技术。
报告不把工具视为彼此孤立的软件，而是将它们组织为“终端完成任务、Markdown 记录过程、\LaTeX{}沉淀正式成果、程序语言解决计算问题”的完整工作流。
通过本次整理，我对研究生阶段如何管理学习、科研、资料和计算环境有了更清晰的认识。

\textbf{关键词：}研究生学习；Linux；Markdown；\LaTeX；Python；C++
\end{summarybox}

%==============目录部分=======================
\begin{multicols}{2}
\small
\tableofcontents
\end{multicols}
\clearpage
\pagenumbering{arabic}

%==================== 第一部分 ====================
\section{第一次课程内容总结与个人认识}

\subsection{研究生阶段的核心任务}

课堂首先比较了本科与研究生阶段的差异。
本科阶段的学习往往以课程、考试和统一安排为主，而研究生阶段更强调自主学习、问题意识和科研产出。
课程数量可能减少，但对阅读、实验、写作和汇报的要求明显提高。
因此，“课少”并不意味着任务少，而是意味着需要自行安排的时间更多。

课堂中最重要的提醒是：\textbf{科研是研究生阶段的核心任务}。
课程成绩需要达到基本要求，但真正体现研究能力的，是能否发现问题、阅读文献、设计方案、完成实验并把结果表达为论文或报告。
科研成果通常也是奖学金评定、毕业审核和后续深造的重要依据。
%=======对比表格=========
\begin{table}[htbp]
\centering
\caption{本科阶段与研究生阶段学习方式对比}
\begin{tabularx}{\textwidth}{C{2.4cm}YY}
\toprule
\textbf{维度} & \textbf{本科阶段} & \textbf{研究生阶段} \\
\midrule
学习驱动 & 课程表、考试与统一作业 & 研究问题、项目任务与个人规划 \\
知识获取 & 教材和课堂讲授为主 & 论文、文档、实验和讨论并重 \\
评价方式 & 课程成绩占比较高 & 科研成果、汇报、论文与综合表现 \\
时间管理 & 安排相对固定 & 自主时间多，更考验计划与执行 \\
最终目标 & 掌握专业基础知识 & 形成独立分析和解决问题的能力 \\
\bottomrule
\end{tabularx}
\end{table}

\subsection{科研进程与阶段性规划}

结合课堂内容，我把这三年的硕士学习概括为以下三步：

\begin{enumerate}
  \item \textbf{第一学年：打基础、定方向。}完成课程学习，熟悉课题组研究方向，阅读高质量论文，掌握必要的实验工具，并在导师指导下确定选题。
  \item \textbf{第二学年：做研究、出成果。}围绕选题开展实验，及时记录数据与失败原因，形成阶段性报告，争取完成论文、专利或学术汇报。
  \item \textbf{第三学年：收尾与衔接。}完成学位论文、送审和答辩，同时结合个人目标准备求职、考公考编或继续深造。
\end{enumerate}

%=======学习流程=======
\begin{center}
\begin{tikzpicture}[node distance=8mm and 6mm,>=Latex,
  stage/.style={rounded corners=2mm,draw=MainBlue,fill=LightBlue,minimum width=4.5cm,minimum height=1.5cm,align=center},
  arrow/.style={->,thick,MainBlue}]
  \node[stage] (a) {第一学年\\课程、文献、研究方向};
  \node[stage,right=of a] (b) {第二学年\\实验、汇报、成果产出};
  \node[stage,below=of b] (c) {第三学年\\论文、答辩、就业衔接};
  \draw[arrow] (a)--(b);
  \draw[arrow] (b)--(c);
\end{tikzpicture}
\end{center}

上述规划不是机械的时间表。
例如，文献阅读与写作应贯穿整个培养过程，而不能等到毕业前才开始。
对我而言，第一学期尤其要做好研究方向调研和工具准备，尽快建立可重复的工作环境。

\subsection{课题组中的协作与沟通}

课堂强调了团结、责任感和规则意识。
课题组不仅是完成个人论文的场所，也是一个协作共同体。按时到实验室、维护公共资源、主动汇报进度、帮助同学解决问题，都会直接影响团队效率。

与导师沟通时，应尽量做到“提前联系、带着问题、说明尝试、记录结论”。有课程时不临时打扰；其他时间需要联系时，应尽量提前预约。
遇到科研问题，可以先查阅文献、文档和既往记录，再把“现象—已做尝试—当前判断—希望获得的帮助”整理清楚。小事先与学长学姐沟通，需要导师决策的事项再正式汇报。

\subsection{计算机管理与正确使用}

课堂第二部分介绍了计算机的日常管理。
我的理解是，研究生电脑不仅是娱乐和办公设备，更是科研数据与实验环境的载体，其稳定性、可恢复性和可复现性都很重要。

\begin{itemize}
  \item \textbf{保持系统可靠：}使用正版或来源可靠的软件，及时更新系统和安全补丁，不随意安装“管家类”捆绑软件。
  \item \textbf{规范卸载与清理：}使用系统卸载功能或可靠卸载工具；用 WizTree 等工具分析空间占用，而不是随意删除未知系统文件。
  \item \textbf{选择合适工具：}文字排版使用 Office 或 \LaTeX，表格和数据处理使用 Excel 或专业工具，尽量让工具与任务匹配。
  \item \textbf{提高操作效率：}熟悉 PowerToys、Windows Terminal等以及常用快捷键；对重复工作尽量脚本化。
  \item \textbf{保护科研资料：}重要论文、代码和实验数据至少保留两份副本；代码使用 Git 记录版本，数据按日期和项目分类。
\end{itemize}

\begin{tipbox}{我的行动计划}
培养建立统一的文件目录与命名规则的习惯；为代码项目使用 Git；定期整理实验记录；每月检查一次备份；把常用命令和排错经验持续补充到个人 Markdown 知识库中。
\end{tipbox}

%==================== 第二部分 ====================
\section{Linux 操作系统与个人速查手册}

\subsection{Linux 的基本认识}

操作系统负责管理处理器、内存、存储设备、文件和外设，并向应用程序提供统一接口。
课堂中学习到常见操作系统大致分为 Windows NT 系列和类 Unix 系统。
Linux 是一种类 Unix 操作系统内核，日常所说的 Ubuntu、Debian、Fedora、Arch Linux 等，实际是由 Linux 内核、GNU 工具、软件包管理器和桌面环境等共同组成的发行版。

Linux 的主要特点包括：多用户、多任务；权限边界清晰；文本工具丰富；适合自动化；服务器生态成熟。
科研中常见的远程服务器、高性能计算平台、容器和深度学习环境大多与 Linux 密切相关。
对初学者而言，Ubuntu 或 Debian 资料丰富、软件管理成熟，适合作为入门选择。这也是后续我们进行科研主要的基础平台。

\subsection{目录结构与路径}

Linux 使用一棵以“/”为根的目录树。理解目录用途，有助于避免误操作。

%======常见的目录结构表==========
\begin{table}[htbp]
\centering
\caption{常见 Linux 目录}
\begin{tabularx}{\textwidth}{C{2.1cm}Y C{2.1cm}Y}
\toprule
\textbf{目录} & \textbf{用途} & \textbf{目录} & \textbf{用途} \\
\midrule
\texttt{/home} & 普通用户主目录 & \texttt{/etc} & 系统与服务配置 \\
\texttt{/usr} & 用户程序、库和共享资源 & \texttt{/var} & 日志、缓存等可变数据 \\
\texttt{/tmp} & 临时文件 & \texttt{/dev} & 设备文件 \\
\texttt{/proc} & 内核和进程的虚拟信息 & \texttt{/mnt} & 临时挂载点 \\
\bottomrule
\end{tabularx}
\end{table}

绝对路径从根目录开始，例如 \texttt{/home/user/project}；相对路径以当前目录为起点。\texttt{.} 表示当前目录，\texttt{..} 表示上一级目录，\texttt{\textasciitilde} 表示当前用户主目录。

\subsection{常用命令速查}

%=========常见的命令表======
\begin{table}[H]
\centering
\caption{个人常用命令速查表}
\small
\begin{tabularx}{\textwidth}{C{3.0cm}Y Y}
\toprule
\textbf{任务} & \textbf{命令示例} & \textbf{说明} \\
\midrule
查看位置 & \texttt{pwd} & 显示当前工作目录 \\
列出文件 & \texttt{ls -lah} & 包括隐藏项，并显示易读大小 \\
切换目录 & \texttt{cd /path/to/dir} & 进入指定目录 \\
创建目录 & \texttt{mkdir -p data/raw} & 递归创建多级目录 \\
复制文件 & \texttt{cp -r src backup/} & 递归复制目录 \\
移动/改名 & \texttt{mv old.txt new.txt} & 移动或重命名 \\
查看文本 & \texttt{less file.log} & 分页查看大文件 \\
搜索文件 & \texttt{find . -name '*.py'} & 按名称查找文件 \\
搜索内容 & \texttt{grep -R 'error' logs/} & 递归查找文本 \\
磁盘空间 & \texttt{df -h}；\texttt{du -sh *} & 查看分区和目录占用 \\
进程信息 & \texttt{ps aux}；\texttt{top} & 查看进程与资源使用 \\
远程登录 & \texttt{ssh user@host} & 通过 SSH 连接服务器 \\
传输文件 & \texttt{scp file user@host:/path} & 在本地与服务器间复制 \\
查看帮助 & \texttt{man command} & 阅读命令手册 \\
\bottomrule
\end{tabularx}
\end{table} \FloatBarrier

%============ TODO 不知道为啥这里位置不对============
\Needspace{8\baselineskip} 管道符 \texttt{|} 可以把前一个命令的输出交给后一个命令处理。例如：

\begin{lstlisting}[language=bash,caption={查找最占空间的五个子目录}]
du -sh ./* | sort -hr | head -n 5
\end{lstlisting}

重定向符 \texttt{>} 覆盖写入文件，\texttt{>>} 追加写入，\texttt{2>} 重定向错误输出。组合命令之前应先在测试目录验证，尤其要谨慎使用带有递归参数的删除命令。

\subsection{权限、软件安装与进程}

执行 \texttt{ls -l} 时，可以看到类似 \texttt{-rwxr-xr--} 的权限字符串。前三位对应所有者，中间三位对应所属组，最后三位对应其他用户；\texttt{r}、\texttt{w}、\texttt{x} 分别表示读、写、执行。

\begin{lstlisting}[language=bash,caption={权限与软件包管理示例}]
# 为脚本增加所有者执行权限
chmod u+x run.sh

# Ubuntu/Debian 更新索引并安装 Git
sudo apt update
sudo apt install git

# 查询占用 8000 端口的进程
ss -lntp | grep ':8000'
\end{lstlisting}

\texttt{sudo} 会以更高权限执行命令，应只在明确知道命令影响范围时使用。软件安装优先使用发行版的软件包管理器，避免从不明网站下载并直接执行脚本。

\subsection{一个可复用的科研工作流}

下面给出我准备采用的 Linux 项目目录与操作流程：

\begin{lstlisting}[language=bash,caption={建立规范的项目目录}]
mkdir -p experiment/{data/raw,data/processed,src,results,docs}
cd experiment
git init
printf "data/raw/\nresults/tmp/\n" > .gitignore
git add .gitignore
git commit -m "init project structure"
\end{lstlisting}

原始数据只读保存，清洗后的数据放入 \texttt{data/processed}，程序放入 \texttt{src}，可复现实验结果放入 \texttt{results}，说明文档和报告放入 \texttt{docs}。这种结构能降低文件混乱和误覆盖的风险。

%==================== 第三部分 ====================
\section{Markdown 学习总结}

\subsection{Markdown、CommonMark 与 GFM}

Markdown 是一种轻量级标记语言，目标是让纯文本既易读，又能转换为 HTML 等格式。
由于早期 Markdown 的边界行为没有完全统一，不同解析器可能产生不同结果。
CommonMark 通过规范化语法与测试用例，减少了解析差异；
GitHub Flavored Markdown（GFM）建立在 CommonMark 基础上，增加了表格、任务列表、删除线和自动链接等 GitHub 常用功能。

%=============对比表格==========
\begin{table}[htbp]
\centering
\caption{CommonMark 与 GFM 的关系}
\begin{tabularx}{\textwidth}{C{2.6cm}YY}
\toprule
\textbf{比较项} & \textbf{CommonMark} & \textbf{GFM} \\
\midrule
定位 & 通用、明确的 Markdown 规范 & 面向 GitHub 场景的扩展规范 \\
核心语法 & 标题、段落、列表、引用、链接、代码等 & 继承 CommonMark 核心语法 \\
典型扩展 & 不统一要求扩展语法 & 表格、任务列表、删除线、自动链接 \\
适用场景 & 跨平台文档与解析器实现 & README、Issue、Pull Request、项目协作 \\
\bottomrule
\end{tabularx}
\end{table}

\subsection{核心语法示例}

\begin{lstlisting}[caption={Markdown 常用语法},language={}]
# 一级标题
## 二级标题

这是**粗体**、*斜体*和`行内代码`。

- 无序列表
1. 有序列表
> 引用内容

[链接文字](https://example.com)
![图片说明](image.png)

| 工具 | 用途 |
| --- | --- |
| Git | 版本控制 |

- [x] 已完成
- [ ] 待完成
\end{lstlisting}

写 Markdown 时，标题符号后应留空格；不同内容块之间保留空行；代码围栏应标明语言。这样既提高源文件可读性，也能获得稳定的渲染结果。

\subsection{扩展功能：数学公式与图表}

部分 Markdown 渲染器支持 \LaTeX{} 数学语法。例如行内公式写作 \verb|$E=mc^2$|，独立公式可以写成：

\begin{lstlisting}[caption={Markdown 中的数学公式},language={}]
$$
\sum_{n=1}^{\infty}\frac{1}{n^2}=\frac{\pi^2}{6}
$$
\end{lstlisting}

其渲染结果为
\[
\sum_{n=1}^{\infty}\frac{1}{n^2}=\frac{\pi^2}{6}.
\]

Mermaid 适合流程图、时序图和状态图；PlantUML 更偏向软件工程中的 UML 图。它们通常通过代码描述图形，便于版本控制和后期修改。例如：

\begin{lstlisting}[caption={Mermaid 流程图源码},language={}]
```mermaid
flowchart LR
    A[提出问题] --> B[查阅文献]
    B --> C[设计实验]
    C --> D[分析结果]
```
\end{lstlisting}

需要注意的是，公式、Mermaid、PlantUML、嵌入 HTML 都属于扩展功能，不同平台的支持程度并不一致。
为了提高可移植性，提交前必须在目标平台预览；面向长期保存的核心内容，应尽量保留纯文本含义。

\subsection{Markdown 的适用边界}

Markdown 特别适合 README、实验日志、会议记录、个人知识库和简单网页内容。
它的优势是输入快、版本差异清晰、无需复杂排版操作；
不足是精细分页、浮动体、交叉引用、复杂公式编号和正式出版排版能力有限。
因此，后续科研中我会用 Markdown 记录过程，用 \LaTeX{}整理正式报告和论文。

%==================== 第四部分 ====================
\section{\LaTeX{} 学习总结与本文实践}

\subsection{基本认识}

\TeX{} 是高质量排版系统，\LaTeX{}在其基础上提供了文档类、宏命令和大量扩展包。它可以理解为一种面向文档结构的“标记语言 + 排版引擎工作流”：作者描述标题、章节、公式、表格和引用的逻辑结构，由编译引擎计算最终版面。

与“看到什么就直接调整什么”的文字处理软件相比，\LaTeX{}更强调内容与格式分离。
其优点是数学公式质量高、编号和交叉引用自动化、长文档样式统一、源文件适合 Git 管理；
缺点是初期需要学习语法，错误信息不够直观，复杂表格或自由版式的调整成本较高。

\subsection{最小中文文档与编译方式}

\begin{lstlisting}[language={[LaTeX]TeX},caption={可编译的最小中文文档（我的第一份\LaTeX{} demo）}]
\documentclass[UTF8]{ctexart}

\title{我的第一份 \LaTeX{} 文档}
\author{任鹏程}
\date{\today}

\begin{document}
\maketitle
你好，世界！
\end{document}
\end{lstlisting}

中文文档建议使用 XeLaTeX 编译。若使用 VS Code，可以安装 LaTeX Workshop 扩展，并让构建配方调用 XeLaTeX；也可以在命令行执行：

\begin{lstlisting}[language=bash,caption={编译命令}]
xelatex report.tex
xelatex report.tex   # 第二次用于更新目录、引用和页码
\end{lstlisting}

\begin{tipbox}{避免重复加载中文支持}
\texttt{ctexart} 本身已经提供中文支持，所以使用 \verb|\documentclass{ctexart}| 后，不需要再写 \verb|\usepackage{ctex}|。二者重复使用可能触发“Package ctex can not be loaded with ctexart”一类错误。
\end{tipbox}

\subsection{文档结构、环境和命令}

\LaTeX{}源文件由导言区和正文区构成。
\verb|\documentclass| 选择文档类型；\verb|\usepackage| 加载功能；\verb|\begin{document}| 与 \verb|\end{document}| 包围正文。
章节使用 \verb|\section| 等命令，列表、公式、表格、图片使用相应环境。

\begin{table}[htbp]
\centering
\caption{常用 \LaTeX{}功能}
\begin{tabularx}{\textwidth}{C{3.3cm}Y Y}
\toprule
\textbf{功能} & \textbf{写法} & \textbf{说明} \\
\midrule
标题信息 & \verb|\title|、\verb|\author| & 需用 \verb|\maketitle| 输出 \\
章节 & \verb|\section{标题}| & 自动编号并进入目录 \\
强调 & \verb|\textbf{内容}| & 语义化粗体强调 \\
行内公式 & \verb|$a^2+b^2=c^2$| & 嵌入段落中的短公式 \\
独立公式 & \verb|\[ ... \]| & 单独居中的无编号公式 \\
引用 & \verb|\label|、\verb|\ref| & 自动维护编号 \\
图片 & \verb|\includegraphics| & 常配合 figure 环境 \\
表格 & tabular、tabularx & 控制列、对齐和宽度 \\
\bottomrule
\end{tabularx}
\end{table}

例如，带编号的公式可以写为
\begin{equation}
  f(x)=\frac{1}{\sqrt{2\pi}\sigma}
  \exp\!\left[-\frac{(x-\mu)^2}{2\sigma^2}\right],
  \label{eq:demo}
\end{equation}
正文使用“式\ref{eq:demo}”引用后，即使中途增加其他公式，编号也会自动更新。

\subsection{本文使用的排版功能}

本报告本身就是一次 \LaTeX{}实践，主要使用了以下功能：

\begin{itemize}
  \item 使用 \texttt{ctexart} 处理中文断行、字号和章节标题；
  \item 使用 \texttt{geometry} 设置 A4 页面边距；
  \item 使用 \texttt{fancyhdr} 生成统一页眉页脚；
  \item 使用 \texttt{tabularx}、\texttt{booktabs} 制作自适应宽度表格；
  \item 使用 \texttt{listings} 展示 Linux、Markdown、Java、JavaScript、Python 与 C++ 代码；
  \item 使用 \texttt{tcolorbox} 制作摘要框和提示框；
  \item 使用 \texttt{hyperref} 生成可点击目录与 PDF 元数据；
  \item 使用 TikZ 绘制阶段规划图，体现代码化制图能力。
\end{itemize}

这些功能说明，\LaTeX{}不仅可以输入数学公式，还能够完成结构化、可复用且风格统一的正式文档。

%==================== 第五部分 ====================
\section{个人编程语言基础与复习总结}

\subsection{已有语言基础概况}

在此前的学习和项目开发中，我已经接触并使用过多种编程语言。其中Java 和 JavaScript 是我较为熟练、实际使用较多的语言；
SQL、HTML 与 CSS 是进行后端数据处理和 Web 开发时经常使用的配套技术。
Python 与 C++ 已有一定基础，本次课程后将继续围绕科研计算和程序设计进行系统复习。

\begin{table}[htbp]
\centering
\caption{个人编程语言与相关技术基础}
\begin{tabularx}{\textwidth}{C{2.5cm}C{2.1cm}Y Y}
\toprule
\textbf{语言或技术} & \textbf{掌握情况} & \textbf{主要用途} & \textbf{相关工具与框架} \\
\midrule
Java & 熟练 & 后端业务开发、接口设计、数据处理 & Spring Boot、MyBatis-Plus、Redis、OkHttp \\
JavaScript & 熟练 & Web 页面交互、前端工程化与浏览器 API & Vue 3、Vite、jQuery、Ant Design Vue \\
SQL & 熟练使用 & 数据查询、统计分析与业务数据维护 & MySQL、复杂查询、分组与关联统计 \\
HTML/CSS & 熟练使用 & 页面结构、布局、响应式界面 & HTML5、CSS3、Less、JSP \\
Python & 有基础 & 自动化、数据分析与科研原型 & NumPy、pandas、Matplotlib \\
C++ & 有基础 & 算法练习与高性能程序 & STL、CMake、编译与调试工具 \\
\bottomrule
\end{tabularx}
\end{table}

\subsection{Java 开发经验总结}

Java 是我实际项目中使用较多的语言。我能够使用面向对象思想组织业务代码，理解封装、继承、多态、集合、异常、泛型、I/O 和并发等常用机制。
在后端开发中，我主要使用 Spring Boot 构建服务，结合 MyBatis-Plus 与 MySQL 完成数据访问，使用 Redis 处理缓存或队列任务，并通过 OkHttp、Gson 等组件调用和解析外部接口。

通过项目实践，我认识到 Java 开发不仅是实现功能，还需要关注分层设计、异常处理、日志记录、事务边界、参数校验和代码可维护性。
例如，控制器负责接收请求，业务逻辑放在服务层，数据访问集中在持久层；对外部接口调用则要设置合理的超时并处理失败情况。
这样的结构有利于测试、排错和后续迭代。

\begin{lstlisting}[language=Java,caption={Java 集合与 Stream API 示例}]
List<Integer> scores = List.of(86, 92, 78, 95);
double average = scores.stream()
        .filter(score -> score >= 80)
        .mapToInt(Integer::intValue)
        .average()
        .orElse(0.0);

System.out.printf("合格数据的平均值：%.2f%n", average);
\end{lstlisting}

\subsection{JavaScript 与 Web 开发经验总结}

JavaScript 主要用于浏览器端交互和前端工程化开发。
我能够处理 DOM、事件、异步请求、JSON 数据和浏览器存储，理解作用域、闭包、原型、模块化、Promise 以及 \texttt{async/await} 等概念。
在传统项目中可以结合 JSP 与 jQuery 完成页面开发，在现代项目中则可以使用 Vue 3、Vite、Ant Design Vue 和 Less 构建组件化界面。
我本人的大部分实际经验也是基于Vue框架并结合ant-design或element-ui进行的项目。

实际开发中，前端代码需要同时考虑用户交互、接口状态和异常情况。
例如，请求开始时显示加载状态，成功后更新页面，失败时给出明确提示，并在 \texttt{finally} 中恢复按钮状态。
相比只追求“页面能运行”，模块划分、数据流向和错误处理更能体现工程质量。

\begin{lstlisting}[language=JavaScript,caption={原生js的 异步请求示例}]
async function loadCandidate(id) {
  try {
    const response = await fetch(`/api/candidates/${id}`);
    if (!response.ok) throw new Error("请求失败");
    return await response.json();
  } catch (error) {
    console.error("加载数据时发生错误：", error);
    throw error;
  }
}
\end{lstlisting}

\subsection{Python 与 C++ 的定位}

Python 语法简洁、生态丰富，适合数据处理、自动化、人工智能、科研原型和可视化；
C++ 更接近底层，能够精细控制内存与性能，适合系统软件、高性能计算、图形开发和对实时性要求较高的场景。
二者并非互相替代：科研中可以用 Python 快速验证思路，在性能瓶颈处使用 C++ 加速。

\begin{table}[htbp]
\centering
\caption{Python 与 C++ 对比}
\begin{tabularx}{\textwidth}{C{2.5cm}YY}
\toprule
\textbf{维度} & \textbf{Python} & \textbf{C++} \\
\midrule
类型系统 & 动态类型，运行时检查较多 & 静态类型，编译期检查较多 \\
开发效率 & 语法简洁，原型开发快 & 语法复杂，需要管理更多细节 \\
执行性能 & 通常较慢，可借助科学计算库 & 原生编译，性能与控制力较强 \\
内存管理 & 自动垃圾回收为主 & RAII 与智能指针，可精细控制 \\
常见用途 & 数据分析、AI、脚本、Web & 系统、高性能、游戏、嵌入式 \\
科研策略 & 组织实验与处理数据 & 加速核心算法或对接底层设备 \\
\bottomrule
\end{tabularx}
\end{table}

\subsection{Python 知识框架}

复习 Python 时，我将知识分为五层：

\begin{enumerate}
  \item \textbf{语法基础：}变量、数值、字符串、列表、字典、集合、条件和循环；
  \item \textbf{程序组织：}函数、模块、包、作用域、异常处理和文件读写；
  \item \textbf{抽象能力：}类、对象、继承、迭代器、生成器和装饰器；
  \item \textbf{工程实践：}虚拟环境、依赖文件、日志、单元测试、类型标注和代码风格；
  \item \textbf{科研生态：}NumPy、pandas、Matplotlib、SciPy 以及具体方向所需框架。
\end{enumerate}

下面的程序演示“读取数据—统计—输出”的基本思路：

\begin{lstlisting}[language=Python,caption={Python 数据统计示例}]
from statistics import mean

scores = {"A": 86, "B": 92, "C": 78, "D": 95}
passed = {name: score for name, score in scores.items()
          if score >= 80}

print(f"平均分：{mean(scores.values()):.2f}")
for name, score in sorted(passed.items(),
                          key=lambda item: item[1], reverse=True):
    print(name, score)
\end{lstlisting}

科研项目中应使用虚拟环境隔离依赖，例如执行 \texttt{python -m venv .venv}，激活后再安装所需包，并用 \texttt{requirements.txt} 或更现代的项目配置记录版本。
原始数据、处理脚本和生成结果应分开保存，保证实验可追溯。

\subsection{C++ 知识框架}

C++ 复习的重点包括基本类型与控制结构、函数、引用与指针、类与对象、模板、标准模板库、异常、文件操作、资源管理和并发。
现代 C++ 更强调 RAII、标准容器、算法和智能指针，应尽量避免手工管理裸指针带来的资源泄漏。

\begin{lstlisting}[language=C++,caption={C++ 标准库算法示例}]
#include <algorithm>
#include <iostream>
#include <numeric>
#include <vector>
#include <stdio.h>

int main() {
    vector<double> values{3.2, 1.5, 4.8, 2.0};
    sort(values.begin(), values.end());
    const double sum = accumulate(values.begin(),
                                       values.end(), 0.0);
    cout << "mean = " << sum / values.size() << '\n';
    return 0;
}
\end{lstlisting}

编译时应开启警告，并明确选择语言标准：

\begin{lstlisting}[language=bash,caption={C++ 编译示例}]
g++ -std=c++17 -Wall -Wextra -O2 main.cpp -o main
./main
\end{lstlisting}

其中 \texttt{-Wall -Wextra} 有助于尽早发现可疑代码，\texttt{-O2} 开启常用优化。
开发过程中可使用 CMake 管理多文件工程，并使用调试器和测试框架定位问题。实际使用中可使用类似Clion等IDE或者VsCode+插件的方式来进行更方便的开发。

\subsection{面向科研的联合工作流}

\begin{center}
\begin{tikzpicture}[node distance=7mm and 8mm,>=Latex,
  box/.style={draw=MainBlue,rounded corners=1.5mm,fill=LightBlue,minimum width=3.6cm,minimum height=1.05cm,align=center},
  arr/.style={->,thick,MainBlue}]
  \node[box] (q) {提出问题\\Markdown 记录};
  \node[box,right=of q] (p) {Python 原型\\数据分析};
  \node[box,below=of p] (c) {C++ 优化\\核心计算};
  \node[box,left=of c] (r) {\LaTeX{} 报告\\公式与结果};
  \draw[arr] (q)--(p);
  \draw[arr] (p)--(c);
  \draw[arr] (c)--(r);
  \draw[arr] (r)--(q);
\end{tikzpicture}
\end{center}

这一流程中，Linux 提供统一运行环境，Git 保存代码版本，Markdown 保存实验日志，Python 完成快速分析，C++处理性能敏感模块，最终由 \LaTeX{}生成正式论文或报告。工具之间形成闭环，才能真正提高科研效率。

%==================== 总结 ====================
\section{总结与后续计划}

通过第一次课程，我认识到研究生阶段的关键不是被动完成安排，而是主动建立问题意识、时间规划和可持续的工作方法。
课堂中关于科研流程、团队沟通和计算机管理的内容，为我之后的学习生活提供了明确提醒。

本次作业进一步把 Linux、Markdown、\LaTeX{}与多种编程语言串联起来。
Java 和 JavaScript 是我已有项目经验的重要基础，可分别承担后端服务与前端交互开发；
SQL、HTML 与 CSS 为完整的 Web 开发提供支撑；Python 与 C++则分别兼顾科研原型的开发效率和核心计算的执行性能。
掌握这些工具的目标并不是追求软件数量，而是根据任务选择合适技术，减少重复劳动、保证实验可复现，并把更多精力放在真正的研究问题上。

下一阶段，我计划按以下节奏持续学习：

\begin{enumerate}
  \item 每周练习一组 Linux 命令，并在虚拟机或服务器中实际操作；
  \item 使用 Markdown 记录课程笔记、论文阅读与实验日志；
  \item 使用 \LaTeX{}完成后续课程报告，逐步学习参考文献、交叉引用和论文模板；
  \item 巩固 Java 与 JavaScript 的工程化能力，同时用 Python 完成数据处理小项目，并复习 C++ 标准库与现代资源管理；
  \item 为所有代码建立 Git 仓库，形成规范的版本记录和 README。
\end{enumerate}

\begin{summarybox}{一句话总结}
把研究任务拆解为可执行计划，把操作过程沉淀为可复用记录，把最终成果组织成可验证、可复现、可交流的文档。
\end{summarybox}

\clearpage
\section*{参考资料}
\addcontentsline{toc}{section}{参考资料}
\begin{enumerate}[label={[\arabic*]}]
  \item CommonMark. \emph{CommonMark Specification}. \url{https://spec.commonmark.org/}
  \item GitHub. \emph{GitHub Flavored Markdown Specification}. \url{https://github.github.com/gfm/}
  \item The Linux man-pages project. \emph{Linux manual pages}. \url{https://www.kernel.org/doc/man-pages/}
  \item The \LaTeX{} Project. \emph{LaTeX documentation}. \url{https://www.latex-project.org/help/documentation/}
  \item Python Software Foundation. \emph{Python Documentation}. \url{https://docs.python.org/3/}
  \item ISO C++ Foundation. \emph{Standard C++ resources}. \url{https://isocpp.org/}
  \item 课程讲义：学业基本情况介绍与计算机基本使用，2026年7月14日。
\end{enumerate}

\end{document}
